	\section{Requisitos}
	\subsection{Requisitos del sistema}
	\noindent El proyecto fue desarrollado y probado en sistema con Linux Debian GNU/Linux 12. La aplicación se encuentra dentro de un contenedor de Docker, con las siguientes versiones:

	\begin{itemize}
		\item docker version 28.2.2, build e6534b4
		\item docker-compose version 1.29.2
	\end{itemize}

	\subsection{Instrucciones de instalaci\'on}
	\begin{enumerate}
		\item Instalar PostgreSQL 17.5.1 desde la \href{https://www.postgresql.org/download/}{p\'agina oficial} de PostgreSQL.
		\item Clonar el repositorio del proyecto desde \href{https://github.com/Ic3manMtz/Servicio-Social.git}{Github} el proyecto se encuentra dentro de la carpeta \texttt{Implementaci\'on}.
		\item Para poder levantar el contenedor de Docker primero es necesario entrar a la carpeta del proyecto, \texttt{cd Implementacion/CLI}.
		\item Una vez dentro de la carpeta hay que darle permisos de ejecución a los archivos, \texttt{start\_container.sh} y \texttt{stop\_container.sh}, con el comando \texttt{chmod +x start\_container.sh stop\_container.sh}.
		\begin{description}
			\item[start\_container.sh] Este script se encarga de crear un contenedor de Docker con la imagen del proyecto y levantarlo.
			\item[stop\_container.sh] Este script se encarga de detener el contenedor de Docker y eliminarlo.
		\end{description}
	\end{enumerate}

\noindent Una vez descargadas las dependencias necesarias. Si es la primera vez que se usará el programa el siguiente paso es crear las tablas de la base de datos utilizando los modelos ya creados.
\begin{enumerate}
	\item Se tiene que abrir el manejar de la base de datos, en el caso de PostgreSQL se llama pgAdmin, para este proyecto se recomienda usar la versión cuatro.
	\item Una vez abierto se tiene que crear una base de datos llamada \texttt{VideoData}. El nombre de la base de datos, el usuario y la contraseña se pueden modificar en la variable de entorno \texttt{DB\_CONNECTION\_STRING}, se encuentra dentro del archivo \texttt{.env}.
	\item El siguiente es ejecutar el archivo \texttt{creationOfModels.py}, se encuentra dentro del paquete \texttt{database} del proyecto.
	\item Al terminar la ejecución ya deben estar las mismas tablas que modelos dentro de la base de datos.
\end{enumerate}
